\documentclass[11pt,a4paper]{article}

% ============================================================
% PACKAGES
% ============================================================
\usepackage[utf8]{inputenc}
\usepackage[T1]{fontenc}
\usepackage[french]{babel}
\usepackage{amsmath, amssymb, amsfonts}
\usepackage[margin=1.3cm, top=1.3cm, bottom=1.1cm]{geometry}
\usepackage{enumitem}
\usepackage{fancyhdr}
\usepackage{tikz}
\usetikzlibrary{arrows.meta}
\usepackage{multicol}
\usepackage{tcolorbox}
\tcbuselibrary{skins}

% ============================================================
% EN-TÊTE / PIED DE PAGE
% ============================================================
\setlength{\headheight}{14pt}
\pagestyle{fancy}
\fancyhf{}
\fancyhead[L]{\small\itshape Évaluation sommative — Équations de droites}
\fancyhead[R]{\small\itshape Première STI2D}
\fancyfoot[C]{\small\thepage}
\renewcommand{\headrulewidth}{0.4pt}

% ============================================================
% LISTES COMPACTES
% ============================================================
\setlist[itemize]{leftmargin=12pt, itemsep=0pt, topsep=1pt, parsep=0pt}
\setlist[enumerate]{leftmargin=16pt, itemsep=2pt, topsep=1pt, parsep=0pt}

\setlength{\parindent}{0pt}

% ============================================================
% DOCUMENT
% ============================================================
\begin{document}

% -------- TITRE --------
\begin{center}
  {\large\bfseries Évaluation sommative — Équations de droites}\\[2pt]
  {\small Durée : \textbf{50 minutes} \quad\textbar\quad
   Première STI2D}
\end{center}

\vspace{-4pt}
\begin{tcolorbox}[colback=gray!5, colframe=gray!40, boxrule=0.4pt,
                  arc=2pt, left=6pt, right=6pt, top=3pt, bottom=3pt]
\small
\textbf{Consignes :} calculatrice interdite · toutes les réponses doivent être
justifiées et rédigées · La réponse finale à chaque question doit être rédigée, commencer par le mot "Finalement" et être entièrement encadrée · Le barème est indicatif (total : \textbf{20 points}) ·
une annexe regroupant les rappels utiles se trouve en fin de sujet.
\end{tcolorbox}
\vspace{4pt}

% ============================================================
% EXERCICE 1 — Lecture graphique
% ============================================================
\section*{Exercice 1 — Lecture graphique \hfill \normalsize (3 points)}

\noindent\small Le plan est muni d'un repère orthonormé. La droite $d$ est représentée
ci-dessous.

\begin{center}
\begin{tikzpicture}[scale=0.55]
  \draw[step=1cm, gray!30, very thin] (-5,-5) grid (5,5);
  \draw[->, thick] (-5.5,0) -- (5.5,0) node[right, font=\small] {$x$};
  \draw[->, thick] (0,-5.5) -- (0,5.5) node[above, font=\small] {$y$};
  \foreach \x in {-5,...,-1,1,...,5}
    \draw (\x,0.1) -- (\x,-0.1) node[below, font=\tiny] {\x};
  \foreach \y in {-5,...,-1,1,...,5}
    \draw (0.1,\y) -- (-0.1,\y) node[left, font=\tiny] {\y};
  \node[below left, font=\tiny] at (0,0) {0};
  \draw[red, thick] (-4,-3) -- (3,4) node[right, font=\small\bfseries] {$d$};
\end{tikzpicture}
\end{center}

\vspace{-4pt}
\begin{enumerate}[label=\textbf{\arabic*.}, leftmargin=18pt]
  \item Lire graphiquement les coordonnées de deux points de la droite $d$.
  \item En déduire le coefficient directeur $m$ de la droite $d$.
  \item Déterminer l'ordonnée à l'origine $p$, puis donner l'équation réduite de $d$.
\end{enumerate}

\vspace{4pt}

% ============================================================
% EXERCICE 2 — Deux points
% ============================================================
\section*{Exercice 2 — Droite passant par deux points \hfill \normalsize (3 points)}

\noindent\small Déterminer l'équation réduite de la droite $(AB)$ dans chacun des cas
suivants.

\begin{enumerate}[label=\textbf{\arabic*.}, leftmargin=18pt]
  \item $A(-1\,;\,3)$ et $B(3\,;\,11)$.
  \item $A(0\,;\,-2)$ et $B(4\,;\,6)$.
\end{enumerate}

\vspace{4pt}

% ============================================================
% EXERCICE 3 — Point et coefficient directeur
% ============================================================
\section*{Exercice 3 — Point et coefficient directeur \hfill \normalsize (3 points)}

\noindent\small Déterminer l'équation réduite de la droite $d$ passant par le point $A$
et ayant pour coefficient directeur $m$.

\begin{enumerate}[label=\textbf{\arabic*.}, leftmargin=18pt]
  \item $A(-2\,;\,4)$ et $m = -2$.
  \item $A(3\,;\,-1)$ et $m = \dfrac{1}{2}$.
\end{enumerate}

\vspace{4pt}

% ============================================================
% EXERCICE 4 — Point et ordonnée à l'origine
% ============================================================
\section*{Exercice 4 — Point et ordonnée à l'origine \hfill \normalsize (2 points)}

\noindent\small Déterminer l'équation réduite de la droite $d$ passant par le point $A$
et dont l'ordonnée à l'origine est $p$.

\begin{enumerate}[label=\textbf{\arabic*.}, leftmargin=18pt]
  \item $A(-1\,;\,5)$ et $p = 2$.
  \item $A(4\,;\,-3)$ et $p = -3$.
\end{enumerate}

\vspace{4pt}

% ============================================================
% EXERCICE 5 — Droites parallèles
% ============================================================
\section*{Exercice 5 — Droite parallèle à une droite donnée \hfill \normalsize (3 points)}

\noindent\small On donne l'équation d'une droite $d$ et un point $A$ n'appartenant pas à $d$.
Déterminer l'équation réduite de la droite $d'$ parallèle à $d$ et passant par $A$.

\begin{enumerate}[label=\textbf{\arabic*.}, leftmargin=18pt]
  \item $d : y = -x + 3$ et $A(2\,;\,5)$.
  \item $d : y = 2x - 1$ et $A(-1\,;\,4)$.
\end{enumerate}

\vspace{4pt}

% ============================================================
% EXERCICE 6 — Tracer une droite
% ============================================================
\section*{Exercice 6 — Tracer une droite dans un repère \hfill \normalsize (3 points)}

\noindent\small Dans le repère ci-dessous, tracer la droite $d$ passant par le point
$A(1\,;\,2)$ et de coefficient directeur $m = -2$. Donner ensuite son équation réduite.

\begin{center}
\begin{tikzpicture}[scale=0.55]
  \draw[step=1cm, gray!30, very thin] (-5,-5) grid (5,5);
  \draw[->, thick] (-5.5,0) -- (5.5,0) node[right, font=\small] {$x$};
  \draw[->, thick] (0,-5.5) -- (0,5.5) node[above, font=\small] {$y$};
  \foreach \x in {-5,...,-1,1,...,5}
    \draw (\x,0.1) -- (\x,-0.1) node[below, font=\tiny] {\x};
  \foreach \y in {-5,...,-1,1,...,5}
    \draw (0.1,\y) -- (-0.1,\y) node[left, font=\tiny] {\y};
  \node[below left, font=\tiny] at (0,0) {0};
  \fill (1,2) circle (2.5pt) node[above right, font=\small\bfseries] {$A$};
\end{tikzpicture}
\end{center}

\vspace{-4pt}
\noindent\small Équation réduite de $d$ : \dotfill

\vspace{4pt}

% ============================================================
% EXERCICE 7 — Problème concret STI2D
% ============================================================
\section*{Exercice 7 — Mise en situation STI2D \hfill \normalsize (3 points)}

\noindent\small Pour permettre l'accès d'un bâtiment aux personnes à mobilité réduite,
un technicien doit concevoir une rampe rectiligne. Dans un repère orthonormé dont
l'unité est le mètre, le point de départ de la rampe est $A(0\,;\,0{,}20)$ et
son point d'arrivée est $B(8\,;\,0{,}60)$.

\begin{enumerate}[label=\textbf{\arabic*.}, leftmargin=18pt]
  \item Déterminer le coefficient directeur $m$ de la droite $(AB)$.
        \emph{Rappel : la pente d'une rampe est égale à son coefficient directeur.}
  \item En déduire l'équation réduite de la droite $(AB)$.
  \item La réglementation impose une pente inférieure ou égale à $5\,\%$
        (soit $m \leq 0{,}05$). La rampe respecte-t-elle cette contrainte ? Justifier.
\end{enumerate}

\vspace{4pt}
\hrule
\vspace{2pt}
\noindent\itshape\small
Rappel : toute démarche cohérente sera valorisée, même si elle n'aboutit pas.


\end{document}